
The same complete transport now varies in $u\ne0$ as well as $t$.
With $\Omega_u=-C_\chi(t)/u^2+B_\chi/u$ from (XD.u1)--(XD.u5),
the precise derivative equations are
\[
 \partial_u\widetilde G=-\Pi^{-*}G\Omega_u\Pi^{-1},\qquad
 \partial_{\bar u}\widetilde G=-\Pi^{-*}\Omega_u^*G\Pi^{-1},
 \quad
 \partial_u\widetilde L_t=\Pi[\Omega_u,L_t]\Pi^{-1},
 \quad
 \partial_u\widetilde A_t=\Pi[\Omega_u,A_t]\Pi^{-1}.
 \tag{LC.u1}
\]
Here $G,A_t,L_t$ are independent of $u$.  Each formula follows
by differentiating its literal definition (LC12), inserting
$\Pi_u=\Pi\Omega_u$ and
$(\Pi^{-1})_u=-\Omega_u\Pi^{-1}$, and retaining both product
terms.  Thus the full numerical-range identity (LC10)--(LC12)
holds at every such $u$ with exactly the same original bound
$E_N(t)$, while its moving metric has the displayed derivatives.
This statement makes no bound on the unweighted target metric.
For the original source lift the exact $u$ commutator is
$r_N[A,\Omega_u]j_E+dK_N^{\rm prim}\Omega_u j_E$
by (DS74); the primitive is in the original cochain degree.

The actual comparison numbers used in (LC14)--(LC16) are now
the eigenvalues $m(u,t),M(u,t)$ of
$G^{-1}\Pi(u,t)^*\Pi(u,t)$.  Their original proof by two
quadratic-form inequalities applies pointwise without alteration.
The full $u$ and mixed curvatures of this same constituent are
(SC.u1)--(SC.u4), with its actual inclusion $I$ and Gram $I^*GI$.
In particular the exact order-two and order-one normal terms
must both be used at $t\ne0$; the original $t$-curvature
comparison (LC15) remains its own calculated matrix entry.
